\chapter{Probability: Foundations and the Law of Large
Numbers}\label{ch:b3:probability}

Year 2 built probability on countable spaces; measure theory
now removes every restriction. A probability space is a
measure space of total mass $1$, random variables are
measurable maps, expectation is the Lebesgue integral --- and
at once the whole analytic arsenal
(\cref{ch:b3:measure,ch:b3:lebesgue,ch:b3:product}) applies to
chance. This chapter installs the dictionary, constructs
infinite sequences of independent random variables (on
$\intcc01$, from binary digits: randomness is hiding inside
Lebesgue measure), proves the Borel--Cantelli lemmas and
Kolmogorov's zero--one law, sorts out the modes of
convergence, and proves the law of large numbers --- the
theorem that makes frequencies converge to probabilities and
statistics possible. The weekend problem gives Etemadi's proof
of the strong law in its definitive $L^1$ form.

\section{The dictionary}

\begin{definition}\label{def:b3:probability:space}
A \emph{probability space}\index{probability space} is a
measure space $(\Omega, \mathcal A, \P)$ with $\P(\Omega) =
1$; elements of $\mathcal A$ are \emph{events}, and a property
holds \emph{almost surely} (a.s.) if its event has
probability $1$. A \emph{random variable}\index{random
variable} is a measurable map $X \colon \Omega \to \R$ (or
$\R^d$: a random vector); its \emph{law}\index{law of a random
variable} is the pushforward probability measure $\P_X =
X_*\P$ on $\R$ (\cref{exo:b3:product:9}), determined by the
\emph{distribution function} $F_X(t) = \P(X \leq t)$
(\cref{exo:b3:measure:3}). $X$ has \emph{density} $f$ if
$\P_X = f\,\dd\lambda$; it is \emph{discrete} if $\P_X$ is a
countable combination of Dirac masses. The
\emph{expectation}\index{expectation} is
\[
\E[X] = \int_\Omega X\,\dd\P
\qquad (X \geq 0 \text{ or } X \in L^1(\P)),
\]
and the \emph{transfer theorem} (\cref{exo:b3:product:9})
computes it in the law: $\E[g(X)] = \int_\R g\,\dd\P_X$ ---
$= \sum g(x_k)p_k$ in the discrete case, $= \int
g(x)f(x)\dd x$ in the density case: Year 2's formulas, now
theorems of one theory. The \emph{variance} is $\V(X) =
\E[(X - \E X)^2] = \E[X^2] - (\E X)^2$ for $X \in L^2$.
\end{definition}

\begin{example}\label{ex:b3:probability:laws}
The standard laws and their transforms of note: Bernoulli
$\mathcal B(p)$, binomial $\mathcal B(n, p)$, geometric,
Poisson $\mathcal P(\lambda)$ (discrete: Year 2's tables
remain valid); uniform on $\intcc01$ (Lebesgue measure
itself); exponential $\mathcal E(\lambda)$ (density
$\lambda\eu^{-\lambda x}\mathbf 1_{x>0}$); the
\emph{Gaussian}\index{Gaussian law} $\mathcal N(m, \sigma^2)$
with density $\frac1{\sigma\sqrt{2\pi}}\exp\bigl(-\frac{(x -
m)^2}{2\sigma^2}\bigr)$ --- a probability density by
\cref{pb:b3:lebesgue:1}, with mean $m$ and variance
$\sigma^2$ (Gaussian moments, \cref{exo:b3:product:10}).
\end{example}

\begin{proposition}[Markov and Chebyshev]
\label{prop:b3:probability:markov}
For $X \geq 0$ and $a > 0$: $\P(X \geq a) \leq \frac{\E
X}{a}$; for $X \in L^2$: $\P\bigl(\abs{X - \E X} \geq
a\bigr) \leq \frac{\V(X)}{a^2}$.\index{Markov's
inequality}\index{Chebyshev's inequality}
\end{proposition}

\begin{proof}
\cref{exo:b3:lebesgue:5}(a); Chebyshev is Markov applied to
$(X - \E X)^2$.
\end{proof}

\section{Independence}

\begin{definition}\label{def:b3:probability:independence}
Sub-$\sigma$-algebras $\mathcal A_1, \dots, \mathcal A_n
\subseteq \mathcal A$ are \emph{independent}\index{independence}
if $\P(A_1\cap\dots\cap A_n) = \prod\P(A_i)$ for all $A_i \in
\mathcal A_i$; events are independent if the $\sigma$-algebras
$\{\varnothing, A_i, A_i^c, \Omega\}$ are; random variables
$X_1, \dots, X_n$ if the $\sigma$-algebras $\sigma(X_i) =
X_i^{-1}(\mathcal B(\R))$ are. An infinite family is
independent if every finite subfamily is.
\end{definition}

\begin{theorem}\label{thm:b3:probability:independence}
$X_1, \dots, X_n$ are independent iff the law of the vector
$(X_1, \dots, X_n)$ is the product measure
$\P_{X_1}\otimes\cdots\otimes\P_{X_n}$. In that case, for
$g_i \geq 0$ (or such that the products are integrable):
\[
\E\Bigl[\prod_ig_i(X_i)\Bigr] = \prod_i\E[g_i(X_i)],
\]
in particular $\E[XY] = \E X\,\E Y$ and $\V(X_1 + \dots +
X_n) = \sum\V(X_i)$ for independent $L^2$ variables.
\end{theorem}

\begin{proof}
If the $X_i$ are independent, the two probability measures
$\P_{(X_1,\dots,X_n)}$ and $\bigotimes\P_{X_i}$ agree on all
products $B_1\times\dots\times B_n$ of Borel sets --- a
$\pi$-system generating $\mathcal B(\R^n)$
(\cref{prop:b3:product:sections}(b)) --- hence everywhere
(\cref{thm:b3:measure:uniqueness}). Conversely, a product law
factorizes all events $\bigcap_iX_i^{-1}(B_i)$: independence. The expectation formula is then
Tonelli/Fubini (\cref{thm:b3:product:tonelli}) through the
transfer theorem; $\E[XY] = \E X\E Y$ is the case $g_i =
\mathrm{id}$, and expanding the square gives the additivity
of variances (cross terms $\E[(X_i - \E X_i)(X_j - \E X_j)]
= 0$).
\end{proof}

\begin{theorem}[Existence of independent sequences]
\label{thm:b3:probability:existence}
On $\bigl(\intcc01, \mathcal L, \lambda\bigr)$ there exists a
sequence $(U_n)_{n\geq1}$ of independent random variables,
each uniform on $\intcc01$. Consequently, for any prescribed
laws $(\mu_n)$ on $\R$ there exist independent $(X_n)$ with
$\P_{X_n} = \mu_n$.\index{independent sequences, existence}
\end{theorem}

\begin{proof}
\emph{Digits.} For $\omega \in \intcc01$, let $(b_k(\omega))$
be its binary digits ($\omega = \sum b_k2^{-k}$; choose the
expansion not ending in all $1$'s --- ambiguity concerns only
a countable, hence null, set). Each $b_k$ is a random
variable ($\{b_k = 1\}$ is a finite union of dyadic
intervals) and the vector $(b_1, \dots, b_m)$ takes each
value in $\{0,1\}^m$ on a dyadic interval of length $2^{-m}$:
the $b_k$ are independent Bernoulli$(\frac12)$.

\emph{Regrouping.} Split $\N^*$ into infinitely many disjoint
infinite sets $(I_n)$ (e.g.\ by prime powers, or diagonals);
let $(k^n_j)_j$ enumerate $I_n$ and set
\[
U_n = \sum_{j\geq1} b_{k^n_j}\,2^{-j} .
\]
Each $U_n$ is uniform: its binary digits are independent fair
bits, so $\P(U_n \in [l2^{-m}, (l+1)2^{-m})) = 2^{-m}$ for
every dyadic interval, and dyadic intervals determine the law
(\cref{thm:b3:measure:uniqueness}). The $U_n$ are independent:
they are functions of disjoint blocks of the independent
family $(b_k)$ --- formally, events $\{U_n \in D_n\}$ for
dyadic $D_n$ depend on finitely many digits from disjoint
sets, and factorize; the $\pi$-system argument upgrades to
all Borel sets.

\emph{Arbitrary laws.} Let $G_n(u) = \inf\{t : F_{\mu_n}(t)
\geq u\}$ (the \emph{quantile function} of the distribution
function $F_{\mu_n}$); the key equivalence $G_n(u) \leq t
\iff u \leq F_{\mu_n}(t)$ (right-continuity of $F$,
monotonicity) shows $X_n = G_n(U_n)$ is measurable with
$\P(X_n \leq t) = \P(U_n \leq F_{\mu_n}(t)) =
F_{\mu_n}(t)$: law $\mu_n$; independence is inherited
(functions of independent variables,
\cref{exo:b3:probability:3}).
\end{proof}

\section{Borel--Cantelli and the zero--one law}

\begin{theorem}[Borel--Cantelli]
\label{thm:b3:probability:borelcantelli}
Let $(A_n)$ be events and $\limsup A_n = \bigcap_N
\bigcup_{n\geq N}A_n$ (``$A_n$ occurs infinitely
often'').\index{Borel--Cantelli lemmas}
\begin{enumerate}
  \item If $\sum\P(A_n) < \infty$, then $\P(\limsup A_n) =
        0$.
  \item If $\sum\P(A_n) = \infty$ \emph{and the $A_n$ are
        independent}, then $\P(\limsup A_n) = 1$.
\end{enumerate}
\end{theorem}

\begin{proof}
(1) is \cref{exo:b3:measure:4}. (2): for $N \leq M$,
independence of complements (\cref{exo:b3:probability:3})
gives
\[
\P\Bigl(\bigcap_{n=N}^{M}A_n^c\Bigr) = \prod_{n=N}^M\bigl(1
- \P(A_n)\bigr) \leq
\exp\Bigl(-\sum_{n=N}^M\P(A_n)\Bigr) \xrightarrow[M \to
\infty]{} 0
\]
($1 - x \leq \eu^{-x}$; the series diverges). So
$\P\bigl(\bigcup_{n\geq N}A_n\bigr) = 1$ for every $N$, and
the decreasing intersection over $N$ still has probability
$1$ (continuity from above, \cref{prop:b3:measure:basics}).
\end{proof}

\begin{theorem}[Kolmogorov's zero--one law]
\label{thm:b3:probability:zeroone}
Let $(X_n)$ be independent and $\mathcal T =
\bigcap_N\sigma(X_N, X_{N+1}, \dots)$ the \emph{tail
$\sigma$-algebra}\index{tail $\sigma$-algebra} (events
insensitive to any finite number of the $X_n$: convergence of
$\sum X_n$, of $\frac{S_n}n$, values of $\limsup$'s, \dots).
Then every $T \in \mathcal T$ has $\P(T) \in \{0,
1\}$.\index{zero--one law}
\end{theorem}

\begin{proof}
Fix $N$. The $\sigma$-algebras $\sigma(X_1, \dots, X_N)$ and
$\sigma(X_{N+1}, \dots)$ are independent: events depending on
disjoint blocks factorize on the generating $\pi$-systems
(cylinders $\bigcap_{i\leq N}\{X_i \in B_i\}$, resp.\ finite
conditions on later variables), and Dynkin
(\cref{thm:b3:measure:dynkin}, applied twice, one side at a
time) extends the factorization. A tail event $T$ lies in
$\sigma(X_{N+1}, \dots)$ for every $N$: $T$ is independent of
every $\sigma(X_1, \dots, X_N)$, hence of the $\sigma$-algebra
they generate, $\sigma(X_1, X_2, \dots)$ (Dynkin once more:
the union of the $\sigma(X_1,\dots,X_N)$ is a $\pi$-system
generating it). But $T \in \sigma(X_1, X_2, \dots)$ too: $T$
is independent \emph{of itself}, $\P(T) = \P(T\cap T) =
\P(T)^2$: $\P(T) \in \{0, 1\}$.
\end{proof}

\section{Modes of convergence}

\begin{definition}\label{def:b3:probability:modes}
$X_n \to X$ \emph{almost surely} if $\P(X_n \to X) = 1$;
\emph{in probability} if $\P(\abs{X_n - X} \geq \varepsilon)
\to 0$ for every $\varepsilon > 0$; \emph{in $L^p$} if
$\E\abs{X_n - X}^p \to 0$.\index{convergence of random
variables}
\end{definition}

\begin{proposition}\label{prop:b3:probability:modes}
(a) a.s.\ convergence implies convergence in probability;
(b) $L^p$ convergence implies convergence in probability;
(c) convergence in probability implies a.s.\ convergence
\emph{along a subsequence};
(d) no other implication holds in general.
\end{proposition}

\begin{proof}
(a) $\P(\abs{X_n - X} \geq \varepsilon) \leq
\P\bigl(\sup_{m\geq n}\abs{X_m - X} \geq \varepsilon\bigr)
\downarrow \P\bigl(\limsup\{\abs{X_m - X} \geq
\varepsilon\}\bigr) = 0$ under a.s.\ convergence (continuity
from above; the limsup event excludes convergence). (b)
Markov: $\P(\abs{X_n - X} \geq \varepsilon) \leq
\varepsilon^{-p}\,\E\abs{X_n - X}^p$. (c) Pick $n_k$ with
$\P(\abs{X_{n_k} - X} \geq 2^{-k}) \leq 2^{-k}$;
Borel--Cantelli (1) makes $\abs{X_{n_k} - X} < 2^{-k}$
eventually, a.s. (d) The typewriter (\cref{exo:b3:lp:3}) on
$(\intcc01, \lambda)$ converges in $L^1$ and in probability
but nowhere pointwise; $n\mathbf 1_{\intoo0{1/n}} \to 0$
a.s.\ but not in $L^1$; details and the remaining
counterexamples in \cref{exo:b3:probability:6}.
\end{proof}

\section{The law of large numbers}

Throughout, $(X_n)$ are independent with the same law
(\emph{i.i.d.}), $S_n = X_1 + \dots + X_n$.

\begin{theorem}[Weak law of large numbers]
\label{thm:b3:probability:wlln}
If $X_1 \in L^2$, with $m = \E X_1$:
\[
\P\Bigl(\Bigl|\frac{S_n}{n} - m\Bigr| \geq
\varepsilon\Bigr) \leq
\frac{\V(X_1)}{n\,\varepsilon^2}
\xrightarrow[n\to\infty]{} 0 :
\]
$\frac{S_n}n \to m$ in probability (and in $L^2$).\index{law
of large numbers (weak)}
\end{theorem}

\begin{proof}
$\E\frac{S_n}n = m$ and $\V\bigl(\frac{S_n}n\bigr) =
\frac{n\V(X_1)}{n^2}$
(\cref{thm:b3:probability:independence}); Chebyshev.
\end{proof}

\begin{theorem}[Strong law of large numbers]
\label{thm:b3:probability:slln}
If $X_1 \in L^1$, then\index{law of large numbers (strong)}
\[
\frac{S_n}{n} \xrightarrow[n\to\infty]{\text{a.s.}} \E[X_1].
\]
We prove it here under the stronger hypothesis $X_1 \in
L^4$; the general case ($L^1$: Etemadi's proof) is the
weekend problem.
\end{theorem}

\begin{proof}[Proof under $\E X_1^4 < \infty$]
Centering ($X_i \mapsto X_i - m$), assume $m = 0$. Expand:
\[
\E[S_n^4] = \sum_{i,j,k,l}\E[X_iX_jX_kX_l]
= n\,\E[X_1^4] + 3n(n-1)\,\bigl(\E[X_1^2]\bigr)^2 \leq
C\,n^2 ,
\]
since independence and centering kill every term containing
an isolated factor ($\E[X_iX_jX_kX_l] =
\E[X_i]\E[\cdots] = 0$ unless the indices pair up: the only
survivors are the $n$ terms $i=j=k=l$ and the $3n(n-1)$
terms with two distinct pairs). Markov:
\[
\P\Bigl(\Bigl|\frac{S_n}n\Bigr| \geq \varepsilon\Bigr)
= \P\bigl(S_n^4 \geq n^4\varepsilon^4\bigr)
\leq \frac{Cn^2}{n^4\varepsilon^4} =
\frac{C}{\varepsilon^4n^2},
\]
summable: Borel--Cantelli (1) gives, for each rational
$\varepsilon$, that $\abs{S_n/n} < \varepsilon$ eventually,
a.s.; intersecting over $\varepsilon \in \Q_+^*$ (countably
many probability-$1$ events): $S_n/n \to 0$ a.s.
\end{proof}

\begin{example}[What the strong law buys]
\label{ex:b3:probability:sllnapps}
(a) \emph{Frequencies}: for i.i.d.\ coin flips, the observed
frequency of heads converges a.s.\ to $p$ --- the empirical
justification of probability itself.
(b) \emph{Monte Carlo}\index{Monte Carlo method}: for $g \in
L^1(\intcc01)$ and $(U_n)$ i.i.d.\ uniform
(\cref{thm:b3:probability:existence}),
$\frac1n\sum_{k\leq n}g(U_k) \to \int_0^1g$ a.s.: integrals
by sampling, in any dimension, at the dimension-independent
rate $\sim n^{-1/2}$ made precise in \cref{ch:b3:clt}.
(c) \emph{Normal numbers}: almost every real number has, in
its binary expansion, asymptotic frequency $\frac12$ of ones
(apply the strong law to the digit variables of
\cref{thm:b3:probability:existence}) --- Borel's theorem, a
statement about \emph{every}day numbers proved by measure:
\cref{pb:b3:probability:1} completes it in all bases.
\end{example}

\begin{method}\label{met:b3:probability:toolkit}
The working order for asymptotic statements about random
sequences: (1) \emph{Is the event a tail event?} Then its
probability is $0$ or $1$
(\cref{thm:b3:probability:zeroone}) and one only has to
decide which. (2) \emph{To prove a.s.\ statements}:
Borel--Cantelli --- summable probabilities for the
``bad'' events, via Markov/Chebyshev-type bounds on
whatever moments exist; independence only needed for the
converse direction. (3) \emph{Subsequence + sandwich}: prove
convergence along a manageable subsequence, control the
oscillation in between by monotonicity or maximal
inequalities --- the skeleton of Etemadi's proof. (4) For
distributional limits, wait for \cref{ch:b3:clt}.
\end{method}

\section{Exercises}

\begin{exercise}[$\star$]\label{exo:b3:probability:1}
(a) Let $X$ have continuous strictly increasing distribution
function $F$. Show that $F(X)$ is uniform on $\intcc01$, and
that $G(U) \sim F$ for $U$ uniform, $G = F^{-1}$: simulation
by inversion.
(b) Compute the distribution function and density of $X^2$
for $X$ uniform on $\intcc{-1}1$, and of $-\frac1\lambda\ln
U$ for $U$ uniform on $\intoo01$.
\end{exercise}

\begin{exercise}[$\star$]\label{exo:b3:probability:2}
(a) Compute mean and variance of the Poisson $\mathcal
P(\lambda)$ and geometric laws via the transfer theorem.
(b) Show that a positive random variable $T$ with $\P(T > t)
> 0$ for all $t$ satisfies the \emph{memoryless} property $\P(T > t + s \mid
T > t) = \P(T > s)$ for all $s, t \geq 0$ iff $T$ is
exponential. \emph{(The survival function satisfies Cauchy's
functional equation; monotonicity replaces continuity.)}
\end{exercise}

\begin{exercise}[$\star\star$]\label{exo:b3:probability:3}
(a) Show that if $X_1, \dots, X_n$ are independent and $f_i$
are Borel functions, the $f_i(X_i)$ are independent.
(b) Show that events $A_1, \dots, A_n$ are independent iff
their complements are, iff the indicators $\mathbf 1_{A_i}$
are independent random variables.
(c) (Pairwise is weaker) Two fair coins: $A = $ first is
heads, $B = $ second is heads, $C = $ the two agree. Show
$A, B, C$ are pairwise independent but not independent.
\end{exercise}

\begin{exercise}[$\star\star$]\label{exo:b3:probability:4}
(a) (Infinite monkey) An i.i.d.\ sequence of uniform
keystrokes on a finite alphabet a.s.\ contains every finite
text infinitely often: prove it with Borel--Cantelli (2) on
disjoint blocks.
(b) (Runs) For i.i.d.\ fair bits, let $R_n$ be the length of
the run of ones starting at position $n$. Show that a.s.\
$R_n \geq (1+\varepsilon)\log_2n$ finitely often, and $R_n
\geq \log_2 n$ infinitely often \emph{(both halves of
Borel--Cantelli; for the second, pass to disjoint blocks to
gain independence)}: the longest run in the first $n$ digits
grows like $\log_2n$.
\end{exercise}

\begin{exercise}[$\star\star$]\label{exo:b3:probability:5}
Let $(X_n)$ be independent.
(a) Show that the radius of convergence of $\sum X_n z^n$ is
an a.s.\ constant (possibly $0$ or $\infty$).
(b) Show that $\P(\sum X_n \text{ converges}) \in \{0, 1\}$
and $\P(S_n/n \to m) \in \{0,1\}$.
(c) Give an event about $(X_n)$ that is \emph{not} a tail
event, and check that the zero--one law can fail for it.
\end{exercise}

\begin{exercise}[$\star\star$]\label{exo:b3:probability:6}
On $(\intcc01, \lambda)$, exhibit --- with proofs --- random
variables such that:
(a) $X_n \to 0$ in probability and in every $L^p$, but
nowhere a.s.;
(b) $X_n \to 0$ a.s.\ but in no $L^p$;
(c) $X_n \to 0$ in $L^1$ but not in $L^2$;
(d) and show: if $X_n \to X$ in probability and $\abs{X_n}
\leq Y \in L^1$, then $X_n \to X$ in $L^1$
\emph{(subsequences + dominated convergence +
the subsubsequence trick)}.
\end{exercise}

\begin{exercise}[$\star\star$]\label{exo:b3:probability:7}
An opinion poll estimates an unknown proportion $p$ by the
empirical frequency $\hat p_n$ of $n$ independent draws.
(a) Chebyshev: show $\P(\abs{\hat p_n - p} \geq \varepsilon)
\leq \frac1{4n\varepsilon^2}$ (use $p(1-p) \leq \frac14$).
(b) How many draws guarantee an error $\leq 3\%$ with
probability $\geq 95\%$ by this bound? (The true answer,
via \cref{ch:b3:clt}, is about $1070$: Chebyshev is honest
but crude.)
\end{exercise}

\begin{exercise}[$\star\star\star$]\label{exo:b3:probability:8}
(Bernstein) For $f \in \mathcal C(\intcc01)$ define the
Bernstein polynomial $B_nf(x) =
\sum_{k=0}^n\binom nkx^k(1-x)^{n-k}f\bigl(\frac kn\bigr)$.
(a) Recognize $B_nf(x) = \E\bigl[f\bigl(\frac
{S_n}n\bigr)\bigr]$ for $S_n$ binomial $\mathcal B(n, x)$.
(b) Prove $B_nf \to f$ \emph{uniformly} on $\intcc01$: split
on $\{\abs{\frac{S_n}n - x} \leq \delta\}$ and its
complement, using uniform continuity and Chebyshev with the
uniform bound $\V(\frac{S_n}n) \leq \frac1{4n}$.
(c) Conclude: a second, probabilistic proof of the
Weierstrass approximation theorem
(\cref{cor:b3:complete:weierstrass}), with the explicit rate
$\norm{B_nf - f}_\infty \leq \frac32\,\omega_f(n^{-1/2})$
for the modulus of continuity $\omega_f$ --- prove at least
the $O(\omega_f(n^{-1/2}))$ form.
\end{exercise}

\begin{exercise}[$\star\star\star$]\label{exo:b3:probability:9}
(Coupon collector) Cards of $n$ types are drawn uniformly
with replacement; let $T_n$ be the number of draws until all
types are seen.
(a) Write $T_n = \sum_{k=1}^{n}\tau_k$ with $\tau_k$
geometric of parameter $\frac{n - k + 1}n$, the $\tau_k$
independent, and deduce $\E T_n = n\,H_n \sim n\ln n$
($H_n$ the harmonic number) and $\V(T_n) \leq
\frac{\pi^2}6n^2$.
(b) Chebyshev: $\frac{T_n}{n\ln n} \to 1$ in probability.
(c) Refine with Borel--Cantelli: show directly $\P(T_n >
\beta n\ln n) \leq n^{1 - \beta}$ for $\beta > 1$
\emph{(union bound on the event that some type is missed
after $\beta n\ln n$ draws, using $1 - x \leq \eu^{-x}$)},
and deduce that along $n = 2^m$, a.s.\ $T_n \leq \beta n\ln
n$ eventually, for every $\beta > 2$.
\end{exercise}

\begin{exercise}[$\star\star$]\label{exo:b3:probability:10}
Using the digit construction
(\cref{thm:b3:probability:existence}):
(a) verify by direct computation that $U = \sum b_{2k}2^{-k}$
(even-indexed digits of a uniform $\omega$) is uniform and
independent of $V = \sum b_{2k-1}2^{-k}$;
(b) deduce a measurable bijection-up-to-null-sets between
$\intcc01$ and $\intcc01^2$ preserving measure, and comment:
one uniform random number contains two (and countably many)
independent ones --- compare with the Peano curve
(\cref{pb:b3:topology:1}), which achieved surjectivity but
not measure-preservation or injectivity.
\end{exercise}

\section{Problem: Etemadi's proof of the strong law}

\begin{problem}[{Weekend problem --- the strong law of large
numbers for i.i.d.\ integrable variables}]
\label{pb:b3:probability:1}
Kolmogorov's strong law --- $\frac{S_n}n \to \E X_1$ a.s.\
for i.i.d.\ $X_n \in L^1$ --- long had only intricate proofs;
in 1981 N.~Etemadi found one of striking economy, using
nothing beyond this chapter (and even weakening independence
to pairwise independence). We follow it. Let $(X_n)$ be
pairwise independent, identically distributed, integrable;
$m = \E X_1$, $S_n = X_1 + \dots + X_n$.

\medskip
\noindent\textbf{Part I --- Reductions.}
\begin{enumerate}
  \item Show that it suffices to treat $X_n \geq 0$
        \emph{(split $X_n = X_n^+ - X_n^-$: check the two
        halves are again pairwise independent i.i.d.\
        integrable)}. Assume henceforth $X_n \geq 0$.
  \item (Truncation) Let $Y_n = X_n\,\mathbf 1_{X_n \leq n}$
        and $S_n^* = Y_1 + \dots + Y_n$. Show
        \[
        \sum_{n\geq1}\P(X_n \neq Y_n) =
        \sum_{n\geq1}\P(X_1 > n) \leq \E[X_1] < \infty
        \]
        (\cref{exo:b3:product:3}), and deduce via
        Borel--Cantelli that $\frac{S_n - S_n^*}{n} \to 0$
        a.s.: it suffices to prove $\frac{S^*_n}n \to m$
        a.s.
  \item Show $\E Y_n = \E\bigl[X_1\mathbf 1_{X_1\leq
        n}\bigr] \to m$ (monotone convergence), hence
        $\frac1n\sum_{k\leq n}\E Y_k \to m$ (Ces\`aro): it
        suffices to prove $\frac{S_n^* - \E S_n^*}{n} \to 0$
        a.s.
\end{enumerate}

\medskip
\noindent\textbf{Part II --- The variance estimate.}
\begin{enumerate}[resume]
  \item Show
        \[
        \V(Y_n) \leq \E[Y_n^2] = \E\bigl[X_1^2\,\mathbf
        1_{X_1 \leq n}\bigr]
        \]
        and, using the layer cake
        (\cref{prop:b3:product:layercake}), the key bound
        \[
        \sum_{n\geq1}\frac{\V(Y_n)}{n^2}
        \leq \sum_{n\geq1}\frac1{n^2}\,
        \E\bigl[X_1^2\mathbf 1_{X_1\leq n}\bigr]
        \leq C\,\E[X_1] < \infty
        \]
        \emph{(exchange the sum and the expectation ---
        Tonelli for series --- and bound $\sum_{n \geq
        x}\frac1{n^2} \leq \frac2{\max(x,1)}$ for the inner
        estimate $x^2\sum_{n\geq x}n^{-2} \leq 2x$)}.
\end{enumerate}

\medskip
\noindent\textbf{Part III --- Convergence along geometric
subsequences.} Fix $\alpha > 1$ and let $k_j =
\lfloor\alpha^j\rfloor$.
\begin{enumerate}[resume]
  \item Using pairwise independence (variances add,
        \cref{thm:b3:probability:independence} --- check that
        additivity of variances needs only pairwise
        independence) and Chebyshev, show for every
        $\varepsilon > 0$:
        \[
        \sum_{j\geq1}\P\Bigl(\Bigl|
        \frac{S^*_{k_j} - \E S^*_{k_j}}{k_j}\Bigr| \geq
        \varepsilon\Bigr)
        \leq
        \frac1{\varepsilon^2}\sum_{j\geq1}\frac1{k_j^2}
        \sum_{n\leq k_j}\V(Y_n)
        = \frac1{\varepsilon^2}\sum_{n\geq1}\V(Y_n)
        \sum_{j\,:\,k_j\geq n}\frac1{k_j^2} .
        \]
  \item Show $\sum_{j : k_j \geq n}k_j^{-2} \leq
        \frac{C_\alpha}{n^2}$ \emph{(geometric series;
        beware the floor: $k_j \geq \frac{\alpha^j}2$ for
        $\alpha^j \geq 2$-type care)}, and conclude with
        question 4 and Borel--Cantelli:
        \[
        \frac{S^*_{k_j} - \E S^*_{k_j}}{k_j}
        \xrightarrow[j\to\infty]{\text{a.s.}} 0,
        \qquad\text{hence}\qquad
        \frac{S^*_{k_j}}{k_j} \to m \ \text{a.s.}
        \]
\end{enumerate}

\medskip
\noindent\textbf{Part IV --- Sandwich and conclusion.}
\begin{enumerate}[resume]
  \item For $k_j \leq n \leq k_{j+1}$, use the monotonicity
        of $S^*_n$ (nonnegative summands!) to show
        \[
        \frac{k_j}{k_{j+1}}\,\frac{S^*_{k_j}}{k_j}
        \;\leq\; \frac{S^*_n}{n} \;\leq\;
        \frac{k_{j+1}}{k_j}\,\frac{S^*_{k_{j+1}}}{k_{j+1}},
        \]
        and deduce, a.s.:
        \[
        \frac m\alpha \leq \liminf\frac{S^*_n}n \leq
        \limsup\frac{S^*_n}n \leq \alpha\,m .
        \]
  \item Let $\alpha \downarrow 1$ along a sequence and
        conclude $\frac{S_n^*}n \to m$ a.s., hence (Part I)
        the \emph{strong law of large numbers}:
        \[
        \boxed{\ \frac{S_n}{n}
        \xrightarrow[n\to\infty]{\text{a.s.}} \E[X_1].\ }
        \]
  \item Where exactly was pairwise independence (rather than
        full independence) sufficient? List the three places
        where independence-type hypotheses were invoked.
\end{enumerate}

\medskip
\noindent\textbf{Part V --- Dividends.}
\begin{enumerate}[resume]
  \item (Borel's normal numbers) Show that
        $\lambda$-almost every $x \in \intcc01$ is
        \emph{normal in every base} $b \geq 2$: each digit
        $0, \dots, b-1$ appears with asymptotic frequency
        $\frac1b$ \emph{(fix $b$ and a digit, apply the
        strong law to the indicator variables --- justify
        that base-$b$ digits of a uniform variable are
        i.i.d.\ uniform on $\{0,\dots,b-1\}$ as in
        \cref{thm:b3:probability:existence} --- then
        intersect the countably many probability-one
        events)}. Exhibit one explicit non-normal number,
        and reflect: the theorem asserts normality of almost
        all numbers, yet proving normality of $\sqrt2$ or
        $\pi$ remains open.
  \item (Monte Carlo, guaranteed) Justify completely the
        method of \cref{ex:b3:probability:sllnapps}(b) for
        $g \in L^1(\intcc01^d)$: construct the i.i.d.\
        uniform sample on $\intcc01^d$ from
        \cref{thm:b3:probability:existence} and
        \cref{exo:b3:probability:10}, and state what the
        strong law delivers.
\end{enumerate}
\end{problem}
